
@@M0@@-algèbres filtrées @@M1@@ en envoyant @@M2@@ sur @@M3@@ et @@M4@@ sur le @@M5@@-ième coefficient de @@M6@@ pour @@M7@@~; utilisant cet homomorphisme, on voit qu'il suffit de vérifier que le morphisme de foncteurs considéré s'annule pour l'élément @@M8@@. On utilise alors l'endomorphisme @@M9@@ de @@M10@@ défini par @@M11@@ pour tout @@M12@@, et le diagramme commutatif~:
@@M13@@
or @@M14@@, @@M15@@ étant l'homomorphisme d'augmentation de @@M16@@~; comme ce dernier annule @@M17@@, il résulte du diagramme analogue écrit pour @@M18@@ que @@M19@@ annule l'image de @@M20@@ dans @@M21@@, d'où le résultat.

Par passage au quotient, on associe donc à tout morphisme de foncteurs~: @@M22@@, un morphisme de foncteurs~: @@M23@@~; par composition avec @@M24@@, et en utilisant l'isomorphisme @@M25@@, on obtient un morphisme~: @@M26@@, ce qui définit l'homomorphisme @@M27@@ cherché. Il est immédiat de vérifier que @@M28@@ et @@M29@@ sont inverses l'un de l'autre.

\textbf{Théorème 1.13.} \uline{Soit} @@M30@@ \uline{un foncteur sur} @@M31@@, \uline{à valeurs dans la catégorie des groupes abéliens filtrés séparés complets}, \uline{et commutant aux invariants sous un groupe fini}. \uline{Alors l'application de} @@M32@@ \uline{dans} @@M33@@ \uline{qui}, \uline{à un morphisme de foncteurs}, \uline{associe sa valeur sur l'élément} 
